% answers_short.tex —— 2005 年·普通高等学校招生考试·北京卷·文科 · 简短答案速查
% 由 tools/make_short_answers.py 自 answers.tex 的 \gksolhead/\gkanswer 提取，
% 并与 2005/2005_数学真题_北京文科_答案要点.md 逐题比对后生成，勿手改（改源 MD 后重跑）。

\begin{gknotice}
\textbf{使用说明}：本节只给最终答案便于快速核对，完整推导见解析版。答案为编写者独立推导，\textbf{非官方标准答案}。
\end{gknotice}

\gkpart{选择题}
\begin{center}\begin{tabular}{ccccccccc}
\toprule
题号 & 1 & 2 & 3 & 4 & 5 & 6 & 7 & 8 \\
\midrule
答案 & C & A & B & C & B & D & C & B \\
\bottomrule
\end{tabular}\end{center}

\gkpart{填空题}
\begin{itemize}[leftmargin=2.6em]
  \item[\textbf{9.}] $x = -1$；$(1, 0)$
  \item[\textbf{10.}] $-20$
  \item[\textbf{11.}] $[-1, 2) \cup (2, +\infty)$
  \item[\textbf{12.}] $\sqrt{2}$
  \item[\textbf{13.}] ②③
  \item[\textbf{14.}] $65$；$20$
\end{itemize}

\gkpart{解答题}
\begin{itemize}[leftmargin=2.6em]
  \item[\textbf{15.}] (1) $-\dfrac{1}{7}$；(2) $\dfrac{7}{6}$
  \item[\textbf{16.}] (1) 见解析版；(2) 见解析版；(3) $\dfrac{2\sqrt{2}}{5}$
  \item[\textbf{17.}] (1) $a_2 = \dfrac{1}{3}$，$a_3 = \dfrac{4}{9}$，$a_4 = \dfrac{16}{27}$；数列通项公式为 $a_n = \begin{cases} 1, & n = 1, \\ \dfrac{1}{3}\left(\dfrac{4}{3}\right)^{n-2}, & n \geqslant 2 \end{cases}$；(2) $\dfrac{3}{7}\left[\left(\dfrac{16}{9}\right)^n - 1\right]$
  \item[\textbf{18.}] (1) $\dfrac{3}{8}$；(2) $\dfrac{20}{27}$；(3) $\dfrac{1}{6}$
  \item[\textbf{19.}] (1) $(-\infty, -1]$ 和 $[3, +\infty)$（或 $(-\infty, -1)$ 和 $(3, +\infty)$）；(2) $-7$
  \item[\textbf{20.}] (1) $W_1$ 表示为 $\begin{cases} x < 0, \\ kx < y < -kx \end{cases}$；$W_2$ 表示为 $\begin{cases} x > 0, \\ -kx < y < kx \end{cases}$；(2) $k^2 x^2 - y^2 = (k^2 + 1)d^2$；(3) 见解析版
\end{itemize}
